\scanpage{6}
factoriser \uline{canoniquement} en
\diagram{
\begin{tikzcd}[ampersand replacement=\&]
(10)\quad X \arrow[r,"f_0","\sim"'] \& Y' \arrow[r,"p"] \& Y
\end{tikzcd}
}
où $Y'$ est un rev. étale fini de $Y$, défini à isom. unique près, et où $X\xrightarrow{\sim}Y'$ est un isom dans la cat. isotopique -- i.e. une classe d'isotopie d'homéomorphismes ? Bien sûr, si on n'a pas \unsure{une d'exigence|une exigence} de canonicité, la factorisabilité (10) de $f$ est triviale -- il suffit de prendre un représentant quelconque de $f$. \deleted{\unsure{Plus|Pour}} Pour préciser la \deleted{question} condition d'unicité, on est amené à se poser la question si pour tt $Y$-autom. \deleted{$\unsure{g}$} $\sigma$ de $Y'$ \deleted{\unsure{la classe}} (\uline{pas} à isot. près !) \deleted{tel que l'on ait} on a l'implicat°
\[
(11)\qquad (\sigma f_0\sim\deleted{\unsure{\sigma}}f_0)\Longrightarrow\sigma=\operatorname{id}_{Y'}\quad ?
\]
Dans cette question, on peut bien sûr supposer $f_0=\operatorname{id}$, donc la question est si un autom d'ordre fini \inserted{$\sigma$} d'\deleted{une} \inserted{la} surface orientée $Y'$ de type fini, qui est isotope à l'identité, est l'identité.

La condition supplémentaire, que $\sigma$ est un $Y$-autom., où $Y'\mathbin{-}Y$ est fini étale, signifie aussi que $\sigma$ opère via une opération \uline{libre} de $\mathbb{Z}/n\mathbb{Z}$ (où $n$ est l'ordre de $\sigma$) -- mais en admettant que si $\sigma$ \inserted{(autom. d'ordre fini)} n'est pas l'identité, ses pts fixes sont isolés,
\margin{5}
